% GENERATED FILE. Edit canonical lessons, then run npm run generate:books.
% canonical-source-hash: fnv1a64:056d038d300b1740
% canonical-lessons: SA-C451-vyathate, SA-C451-asvasiti, SA-C451-anunayati, SA-C451-vyakulah, SA-C451-santapah

\chapter{Distress and Comfort}
\label{ch:sa-distress-and-comfort}
\hlchaptermodality{451}

\begin{chapteropening}
\textbf{By the end of this chapter:} \emph{I can say how someone was comforted.}

Complete the last lesson of chapter 451: I can say how someone was comforted.
\end{chapteropening}

\section[\texorpdfstring{\sk{व्यथते} (vyathate)}{व्यथते (vyathate)}]{\sk{व्यथते} (vyathate) — is pained; is distressed}

\label{lesson:SA-C451-vyathate}



\begin{quote}
Before the new one: say the Sanskrit for tears, then the Sanskrit for a sigh.
\end{quote}

\subsection*{You'll want to know: \sk{व्यथते}}
\textbf{\sk{व्यथते}} — \emph{vyathate} — \textquotedblleft{}is pained; is distressed\textquotedblright{}.

\begin{quote}
\sk{मित्रस्य} \sk{दुःखेन} \sk{मीरा} \sk{व्यथते।}
\end{quote}

\begin{quote}
\emph{mitrasya duḥkhena mīrā vyathate.}
\end{quote}

\begin{quote}
Mira is pained by her friend's sorrow.
\end{quote}

Past (she …): \sk{सा} \sk{व्यथितवती}, \emph{sā vyathitavatī}.

\begin{cousinweb}
\textbf{\sk{व्यथते}} (\emph{vyathate}) comes from the root √vyath, \textquotedblleft{}to tremble\textquotedblright{}; MW's ppp \emph{vyathita} is \textquotedblleft{}disquieted, distressed, afflicted\textquotedblright{} (MW).
\end{cousinweb}

\begin{grammarlens}[title={What you've built}]
One more word for describing what you lived and felt.
\end{grammarlens}

\subsection*{Your turn}
\begin{itemize}
\raggedright
  \item \emph{Say it:} \emph{vyathate}
  \item \emph{Say it:} \emph{vyathate}, once more
  \item \emph{Say it:} \emph{mitrasya duḥkhena mīrā vyathate.} — Mira is pained by her friend's sorrow.
  \item \emph{Recall:} say the Sanskrit for tears, then the Sanskrit for a sigh, then say \emph{vyathate} again
\end{itemize}

\begin{tcolorbox}[breakable,colback=teal!4,colframe=teal!35!black,title={Before you move on}]
What does \sk{व्यथते} mean? (\textquotedblleft{}Is pained; is distressed\textquotedblright{}.) Say it once more.
\end{tcolorbox}
\section[\texorpdfstring{\sk{आश्वसिति} (āśvasiti)}{आश्वसिति (āśvasiti)}]{\sk{आश्वसिति} (āśvasiti) — takes heart}

\label{lesson:SA-C451-asvasiti}



\begin{quote}
Before the new one: say the Sanskrit for a sigh, then the Sanskrit for is pained.
\end{quote}

\subsection*{You'll want to know: \sk{आश्वसिति}}
\textbf{\sk{आश्वसिति}} — \emph{āśvasiti} — \textquotedblleft{}takes heart\textquotedblright{}.

\begin{quote}
\sk{वैद्यः} \sk{आगच्छति}, \sk{माता} \sk{आश्वसिति} \sk{च।}
\end{quote}

\begin{quote}
\emph{vaidyaḥ āgacchati, mātā āśvasiti ca.}
\end{quote}

\begin{quote}
The physician comes, and the mother takes heart.
\end{quote}

Past (she …): \sk{सा} \sk{आश्वस्तवती}, \emph{sā āśvastavatī}.

\begin{cousinweb}
\textbf{\sk{आश्वसिति}} (\emph{āśvasiti}) is built from \emph{ā}, \textquotedblleft{}towards\textquotedblright{}, and the root √śvas, \textquotedblleft{}to breathe\textquotedblright{}: to breathe again; MW gives \textquotedblleft{}to take or recover breath\textquotedblright{} and \textquotedblleft{}take heart\textquotedblright{} (MW).
\end{cousinweb}

\begin{grammarlens}[title={What you've built}]
One more word for describing what you lived and felt.
\end{grammarlens}

\subsection*{Your turn}
\begin{itemize}
\raggedright
  \item \emph{Say it:} \emph{āśvasiti}
  \item \emph{Say it:} \emph{āśvasiti}, once more
  \item \emph{Say it:} \emph{vaidyaḥ āgacchati, mātā āśvasiti ca.} — The physician comes, and the mother takes heart.
  \item \emph{Recall:} say the Sanskrit for a sigh, then the Sanskrit for is pained, then say \emph{āśvasiti} again
\end{itemize}

\begin{tcolorbox}[breakable,colback=teal!4,colframe=teal!35!black,title={Before you move on}]
What does \sk{आश्वसिति} mean? (\textquotedblleft{}Takes heart\textquotedblright{}.) Say it once more.
\end{tcolorbox}
\section[\texorpdfstring{\sk{अनुनयति} (anunayati)}{अनुनयति (anunayati)}]{\sk{अनुनयति} (anunayati) — coaxes; soothes}

\label{lesson:SA-C451-anunayati}



\begin{quote}
Before the new one: say the Sanskrit for is pained, then the Sanskrit for takes heart.
\end{quote}

\subsection*{You'll want to know: \sk{अनुनयति}}
\textbf{\sk{अनुनयति}} — \emph{anunayati} — \textquotedblleft{}coaxes; soothes\textquotedblright{}.

\begin{quote}
\sk{शिशुः} \sk{रोदिति}, \sk{माता} \sk{तम्} \sk{अनुनयति।}
\end{quote}

\begin{quote}
\emph{śiśuḥ roditi, mātā tam anunayati.}
\end{quote}

\begin{quote}
The baby cries, and the mother soothes it.
\end{quote}

Past (she …): \sk{सा} \sk{शिशुम्} \sk{अनुनीतवती}, \emph{sā śiśum anunītavatī}.

\begin{cousinweb}
\textbf{\sk{अनुनयति}} (\emph{anunayati}) is built from \emph{anu}, \textquotedblleft{}after\textquotedblright{}, and the root √nī, \textquotedblleft{}to lead\textquotedblright{}, the root of \textbf{\sk{नयति}}; MW gives \textquotedblleft{}to conciliate, pacify, win over\textquotedblright{} (MW).
\end{cousinweb}

\begin{grammarlens}[title={What you've built}]
One more word for describing what you lived and felt.
\end{grammarlens}

\subsection*{Your turn}
\begin{itemize}
\raggedright
  \item \emph{Say it:} \emph{anunayati}
  \item \emph{Say it:} \emph{anunayati}, once more
  \item \emph{Say it:} \emph{śiśuḥ roditi, mātā tam anunayati.} — The baby cries, and the mother soothes it.
  \item \emph{Recall:} say the Sanskrit for is pained, then the Sanskrit for takes heart, then say \emph{anunayati} again
\end{itemize}

\begin{tcolorbox}[breakable,colback=teal!4,colframe=teal!35!black,title={Before you move on}]
What does \sk{अनुनयति} mean? (\textquotedblleft{}Coaxes; soothes\textquotedblright{}.) Say it once more.
\end{tcolorbox}
\section[\texorpdfstring{\sk{व्याकुलः} (vyākulaḥ)}{व्याकुलः (vyākulaḥ)}]{\sk{व्याकुलः} (vyākulaḥ) — troubled; perplexed}

\label{lesson:SA-C451-vyakulah}



\begin{quote}
Before the new one: say the Sanskrit for takes heart, then the Sanskrit for coaxes.
\end{quote}

\subsection*{You'll want to know: \sk{व्याकुलः}}
\textbf{\sk{व्याकुलः}} — \emph{vyākulaḥ} — \textquotedblleft{}troubled; perplexed\textquotedblright{}.

\begin{quote}
\sk{स्यूतः} \sk{नष्टः}, \sk{अतः} \sk{यात्री} \sk{व्याकुलः} \sk{अस्ति।}
\end{quote}

\begin{quote}
\emph{syūtaḥ naṣṭaḥ, ataḥ yātrī vyākulaḥ asti.}
\end{quote}

\begin{quote}
The bag is lost, so the passenger is troubled.
\end{quote}

\begin{cousinweb}
\textbf{\sk{व्याकुलः}} (\emph{vyākulaḥ}) is built from \emph{vi} and \emph{ākula}, \textquotedblleft{}crowded, confused\textquotedblright{}; MW gives \textquotedblleft{}bewildered, perplexed, troubled\textquotedblright{} (MW).
\end{cousinweb}

\begin{grammarlens}[title={What you've built}]
One more word for describing what you lived and felt.
\end{grammarlens}

\subsection*{Your turn}
\begin{itemize}
\raggedright
  \item \emph{Say it:} \emph{vyākulaḥ}
  \item \emph{Say it:} \emph{vyākulaḥ}, once more
  \item \emph{Say it:} \emph{syūtaḥ naṣṭaḥ, ataḥ yātrī vyākulaḥ asti.} — The bag is lost, so the passenger is troubled.
  \item \emph{Recall:} say the Sanskrit for takes heart, then the Sanskrit for coaxes, then say \emph{vyākulaḥ} again
\end{itemize}

\begin{tcolorbox}[breakable,colback=teal!4,colframe=teal!35!black,title={Before you move on}]
What does \sk{व्याकुलः} mean? (\textquotedblleft{}Troubled; perplexed\textquotedblright{}.) Say it once more.
\end{tcolorbox}
\section[\texorpdfstring{\sk{सन्तापः} (santāpaḥ)}{सन्तापः (santāpaḥ)}]{\sk{सन्तापः} (santāpaḥ) — anguish}

\label{lesson:SA-C451-santapah}



\begin{quote}
Before the new one: say the Sanskrit for coaxes, then the Sanskrit for troubled.
\end{quote}

\subsection*{You'll want to know: \sk{सन्तापः}}
\textbf{\sk{सन्तापः}} — \emph{santāpaḥ} — \textquotedblleft{}anguish\textquotedblright{}.

\begin{quote}
\sk{मम} \sk{हृदये} \sk{महान्} \sk{सन्तापः} \sk{अस्ति।}
\end{quote}

\begin{quote}
\emph{mama hṛdaye mahān santāpaḥ asti.}
\end{quote}

\begin{quote}
There is great anguish in my heart.
\end{quote}

A masculine noun.

\begin{cousinweb}
\textbf{\sk{सन्तापः}} (\emph{santāpaḥ}) is built from \emph{sam} and the root √tap, \textquotedblleft{}to burn\textquotedblright{}, the root of \textbf{\sk{तापः}}: a burning through and through; MW gives \textquotedblleft{}affliction, sorrow, anguish\textquotedblright{} (MW).
\end{cousinweb}

\begin{grammarlens}[title={What you've built}]
One more word for describing what you lived and felt.
\end{grammarlens}

\subsection*{Your turn}
\begin{itemize}
\raggedright
  \item \emph{Say it:} \emph{santāpaḥ}
  \item \emph{Say it:} \emph{santāpaḥ}, once more
  \item \emph{Say it:} \emph{mama hṛdaye mahān santāpaḥ asti.} — There is great anguish in my heart.
  \item \emph{Recall:} say the Sanskrit for coaxes, then the Sanskrit for troubled, then say \emph{santāpaḥ} again
\end{itemize}

\begin{tcolorbox}[breakable,colback=teal!4,colframe=teal!35!black,title={Before you move on}]
What does \sk{सन्तापः} mean? (\textquotedblleft{}Anguish\textquotedblright{}.) Say it once more.
\end{tcolorbox}
